% Exercises the label of a citation node that carries a note: the note is set
% inside the tag and after it ("[1, pp. 23-27]"), the way a citation with a note
% is printed, and its -- and --- ligatures become the HTML character entities
% Graphviz decodes, so the graph shows the dashes the document shows. Two cites
% of one work with different notes must still give two nodes, and a note-less
% cite of the same work a third. A manual \bibitem keeps it bibtex-free.
\documentclass{article}
\usepackage{amsthm}
\usepackage[cite,citelabel=tag]{proofgraph}
\newtheorem{theorem}{Theorem}
\begin{document}
\begin{theorem}\label{t}A theorem.\end{theorem}
\begin{proof}
  See \cite[pp.~23--27]{knuth}, \cite[Chapters 1---3]{knuth}
  and \cite{knuth}.
\end{proof}
\begin{thebibliography}{9}\bibitem{knuth}A reference.\end{thebibliography}
\end{document}
